\chapter{Universal ORIUS Across Domains}
\label{ch:universal-orius-across-domains}

\section{What universality means in this dissertation}
ORIUS is universal in a specific sense. It does not assume that batteries, vehicles,
industrial plants, healthcare systems, navigation stacks, and aerospace controllers share a
plant model, control law, or optimization target. It assumes only that each domain must
answer the same release-boundary question once observation becomes unreliable: what action
may still be released safely, under what margin, with what fallback, and under what
certificate? The shared object is therefore not the nominal controller. It is the runtime
safety logic that mediates degraded observation before actuation.

\section{The shared domain template}
Every domain row in this dissertation is described through the same template:
\begin{enumerate}
  \item the local degraded-observation hazard;
  \item the ORIUS instantiation wrapped around nominal intelligence;
  \item the defended safety object;
  \item the current evidence status; and
  \item the universality contribution that row makes to the thesis.
\end{enumerate}
That template is what makes cross-domain comparison scientific rather than rhetorical. The
rows are not interchangeable, but they are comparable because ORIUS forces them to answer
the same runtime question through one implementation grammar: a \texttt{DomainAdapter}
binding, one Detect--Calibrate--Constrain--Shield--Certify step, one benchmark schema, and
one certificate-aware release semantics.

\section{Cross-domain status at a glance}
\begin{table}[htbp]
\centering
\scriptsize
\caption{Equal-domain parity gate for the universal-first ORIUS book.}
\label{tab:orius-equal-domain-parity}
\resizebox{\textwidth}{!}{%
\begin{tabular}{p{2.5cm}ccccccccp{3.5cm}}
\toprule
\textbf{Domain} & \textbf{Adapter} & \textbf{Raw/Data} & \textbf{Train} & \textbf{Replay} & \textbf{Soundness} & \textbf{Fallback} & \textbf{CertOS} & \textbf{P5} & \textbf{Tier / Exact blocker}\\
\midrule
Battery Energy Storage & pass & locked\_reference & pass & pass & reference\_witness & safe\_hold\_validated & evaluated & evaluated & reference; battery\_reference\_row\\
Autonomous Vehicles & pass & verified & pass & pass & pass\_under\_ttc\_entry\_barrier\_contract & bounded\_brake\_fallback\_validated & evaluated & gated\_pending\_shared\_constraint\_surface & proof\_validated; multi\_lane\_and\_higher\_dimensional\_repair\_open\\
Industrial Process Control & pass & verified & pass & pass & pass & bounded\_runtime\_pass & evaluated & evaluated & proof\_validated; bounded\_to\_current\_plant\_family\\
Medical and Healthcare Monitoring & pass & verified & pass & pass & pass & bounded\_runtime\_pass & evaluated & gated\_pending\_shared\_constraint\_surface & proof\_validated; bounded\_to\_current\_monitoring\_and\_intervention\_contract\\
Navigation and Guidance & pass & blocked\_real\_data\_gap & blocked\_real\_data\_gap & portability\_only & blocked\_real\_data\_gap & bounded\_runtime\_pass & gated & gated & shadow\_synthetic; navigation\_real\_data\_row\_missing\\
Aerospace Control & pass & placeholder\_surface & placeholder\_surface & experimental\_replay\_only & experimental\_placeholder & bounded\_runtime\_pass & gated & gated & experimental; real\_multi\_flight\_safety\_task\_missing\\
\bottomrule
\end{tabular}}
\end{table}


The parity matrix summarizes the current empirical boundary of the dissertation. Battery is
the witness row. Industrial process control and healthcare monitoring are defended bounded
rows. The autonomous-vehicle row is a defended bounded runtime row under the current
artifact surface, while canonical raw-data closure remains incomplete. Navigation remains
shadow-synthetic, and aerospace remains experimental. The matrix should therefore be read as
the evidence boundary for a universal runtime architecture, not as proof that every row has
already reached equal maturity.

\section{Domain synthesis}
\paragraph{Battery energy storage.}
The battery hazard is hidden constraint violation under degraded feeder, load, or state
telemetry. ORIUS instantiates the runtime with tight state margins, explicit repair, and
auditable release over a heavily instrumented control surface. Its evidence status is the
deepest in the dissertation, which is why battery serves as the witness row. Its
universality contribution is to show that the full theorem-to-code-to-artifact chain can be
closed on a real physical control problem.

\paragraph{Autonomous vehicles.}
The autonomous-vehicle hazard is unsafe longitudinal release under stale or degraded scene
understanding. ORIUS instantiates the runtime through bounded time-to-collision and
predictive-entry-barrier logic, with explicit brake-oriented fallback and certificate-aware
release checks. Its evidence status is a defended bounded runtime row under the current
artifact surface, while canonical raw-data closure remains incomplete. Its universality
contribution is to show that the same release semantics remain meaningful in a fast-horizon,
continuous-action mobility setting.

\paragraph{Industrial process control.}
The industrial hazard is constraint excursion under corrupted or delayed plant telemetry.
ORIUS instantiates the runtime through bounded power or process envelopes, repair against
shared operating constraints, and governed fallback. Its evidence status is a defended
bounded row. Its universality contribution is to show that ORIUS can operate in tightly
regulated process loops where plant constraints are explicit and intervention must remain
operationally interpretable.

\paragraph{Healthcare monitoring.}
The healthcare hazard is clinically unsafe release under partially observed physiological
state, especially when alerts and interventions are triggered from noisy or degraded
streams. ORIUS instantiates the runtime through physiological margin checks, alert fallback,
and certificate-aware release semantics. Its evidence status is a defended bounded row. Its
universality contribution is to show that the ORIUS grammar survives in a setting where the
protected object is patient-facing monitoring behavior rather than actuation on an energy or
motion platform.

\paragraph{Navigation and guidance.}
The navigation hazard is trajectory release under degraded localization or scene reference.
ORIUS instantiates the runtime through hold-position or corridor-preserving fallback logic.
Its current evidence status remains shadow-synthetic because the defended real-data
train--replay chain is not yet closed. Its universality contribution is still valuable: it
shows how the ORIUS release boundary maps onto spatial feasibility problems even before the
row is fully defended.

\paragraph{Aerospace control.}
The aerospace hazard is unsafe envelope or guidance release under degraded flight
observation. ORIUS instantiates the runtime through envelope-aware repair, bounded fallback,
and certificate semantics tied to flight safety objects. Its current evidence status remains
experimental because a defended multi-flight runtime validation surface is still missing.
Its universality contribution is to show that ORIUS extends to high-consequence control
domains without pretending that structural portability already implies empirical closure.

\begin{figure}[htbp]
\centering
\includegraphics[width=0.92\textwidth]{fig_multi_domain_validation}
\caption{Cross-domain ORIUS validation surface. One runtime grammar carries across six
domains, while each row remains limited by its own defended evidence surface.}
\label{fig:multi-domain-validation}
\end{figure}

The figure should be read as an interpretation aid rather than as a replacement for the
parity table. It makes visible that the same runtime object appears in every domain, even
though the domains differ sharply in dynamics, sensing, and safety geometry.

\section{Cross-domain interpretation}
Two conclusions follow from the domain synthesis. First, what transfers is not a controller
family but a release-boundary grammar: reliability loss, uncertainty inflation, action
repair, fallback semantics, and certificate-aware release. Second, what remains local is the
geometry of the safety object itself: a battery margin is not a vehicle headway constraint,
and a vehicle headway constraint is not a physiological alert threshold or a flight
envelope. ORIUS matters scientifically because it preserves one runtime meaning while
allowing those local geometries to remain local.

The most important implication is that universality here is structural rather than
homogenizing. ORIUS does not make domains identical. It makes them legible under one safety
question: when reality and observation diverge, what action may still be released, and with
what certificate? That is why a battery witness row, a bounded vehicle row, industrial and
healthcare rows, and still-gated navigation and aerospace rows can all belong in one
dissertation without collapsing into a single benchmark story.

\section{Blocked rows and the present thesis boundary}
The blocked rows are simple to state and should remain visible. Navigation is not yet a
defended peer because the real-data closure row is still missing. Aerospace is not yet a
defended peer because the dissertation does not yet contain a defended multi-flight runtime
surface with material post-repair gain. That boundary does not weaken the universal thesis.
It clarifies it: ORIUS already establishes a reusable runtime safety grammar across domains,
while equal-domain closure remains an empirical task rather than a rhetorical shortcut.
